% Présentation Beamer : problème -> architecture -> méthode -> application -> résultats -> conclusion.
% Compilation : make slides. Démonstration obligatoire (diapositives 9 et 10). ~15 diapositives.
\documentclass[aspectratio=169,11pt]{beamer}
\usepackage[utf8]{inputenc}
\usepackage[T1]{fontenc}
\usepackage[french]{babel}
\usepackage{amsmath,amssymb,booktabs,graphicx,xcolor}
\usetheme{Madrid}
\usecolortheme{whale}
\setbeamertemplate{navigation symbols}{}
\graphicspath{{../artifacts/figures/}{../report/figures/}}

\newcommand{\R}{\mathbb{R}}
\newcommand{\norm}[1]{\lVert #1\rVert_2}
\DeclareMathOperator*{\argmax}{arg\,max}
\newcommand{\TODO}[1]{\textcolor{red}{[#1]}}
% \figopt{fichier}{hauteur}{} : inclut la figure si elle existe
\newcommand{\figopt}[2]{\IfFileExists{../artifacts/figures/#1}{\includegraphics[height=#2]{#1}}%
  {\IfFileExists{../report/figures/#1}{\includegraphics[height=#2]{#1}}{\TODO{figure \detokenize{#1} absente}}}}
\newcommand{\inputopt}[1]{\IfFileExists{#1}{\input{#1}}{\TODO{\detokenize{#1} absent}}}

\title[Moteur multimodal]{Moteur multimodal de recherche et d'assistance}
\subtitle{Texte, image et audio dans un espace commun}
\author[\'Equipe Projet 3]{CLAUDY \and IDA \and FRESNEL \and RAPHA\"EL}
\institute[ENSPY]{ENSPY -- INF50101 -- Projet 3}
\date{2 d\'ecembre 2026}

\begin{document}

% 1 ----------------------------------------------------------------- Orateur : FRESNEL
\begin{frame}\titlepage\end{frame}

\begin{frame}{Plan}
  \tableofcontents
\end{frame}

\section{Problème et données}
% 2 ------------------------------------------------------------------ Orateur : IDA
\begin{frame}{Problème : chercher dans un corpus multimédia}
  \begin{itemize}
    \item Représenter \textbf{texte, image et audio} dans un \textbf{même espace} pour rechercher, comparer, interroger.
    \item Tâches : texte$\to$image, image$\to$texte, audio$\to$contenu ; tâche secondaire : classification audio zero-shot.
    \item Livrables : application Flask, évaluation Recall@K / Precision@K, comparaison mono-modal contre multimodal.
  \end{itemize}
  \TODO{1 phrase : cas d'usage concret}
\end{frame}

% 3 ------------------------------------------------------------------ Orateur : IDA
\begin{frame}{Données : le texte comme pivot}
  \begin{center}\small
  \begin{tabular}{llll}
    \toprule
    Jeu & Rôle & Taille & Limites\\
    \midrule
    Flickr8k & image$\leftrightarrow$texte & 8\,091 images, 5 légendes & anglais, biais occidental\\
    Clotho v2 & audio$\leftrightarrow$texte & $\approx$ 6\,000 clips, 5 légendes & audio Freesound\\
    ESC-50 & zero-shot & 2\,000 clips, 50 classes & recouvrement Freesound\\
    \bottomrule
  \end{tabular}
  \end{center}
  Aucun jeu n'aligne les trois modalités : image$\leftrightarrow$audio se déduit \textbf{par le texte}.
\end{frame}

\section{Architecture}
% 4 ---------------------------------------------------------------- Orateur : CLAUDY
\begin{frame}{Architecture : hors-ligne et en ligne}
  \centering\figopt{architecture.png}{0.78\textheight}
\end{frame}

% 5 ---------------------------------------------------------------- Orateur : CLAUDY
\begin{frame}{Index et recherche}
  \begin{itemize}
    \item 4 index FAISS exacts \texttt{IndexFlatIP} : images, audio, légendes, items fusionnés ($\approx 202$\,Mo).
    \item Vecteurs normalisés : $\norm{z_a-z_b}^2=2-2\,z_a\!\cdot\!z_b$, donc cosinus = produit scalaire = même classement.
    \item Recherche exacte : $\approx 14$\,ms sur 70\,455 légendes (limitée par la mémoire) ; filtres par pré-filtrage exact.
    \item Audio : AST $=1{,}30\cdot10^{11}$ MACs pour 10 s, soit $\approx 5$\,s sur CPU, donc collection encodée hors-ligne.
  \end{itemize}
\end{frame}

\section{Méthode}
% 6 --------------------------------------------------------------- Orateur : FRESNEL
\begin{frame}{Encodeurs et espace commun ($d=512$)}
  \begin{center}\small
  \begin{tabular}{lll}
    \toprule
    Modalité & Encodeur (gelé) & Sortie\\
    \midrule
    Texte & CLIP ViT-B/32, tour texte (ancre) & 512\\
    Image & CLIP ViT-B/32, tour visuel & 512\\
    Audio & AST + tête MLP entraînée & 768 $\to$ 512\\
    Baseline texte & SBERT all-MiniLM-L6-v2 & 384 (espace séparé)\\
    \bottomrule
  \end{tabular}
  \end{center}
  Normalisation L2 : $z=u/\norm{u}$.
\end{frame}

% 7 --------------------------------------------------------------- Orateur : FRESNEL
\begin{frame}{Fusion : tardive, précoce, attention croisée}
  \small
  \begin{tabular}{@{}lccc@{}}
    \toprule
    & Précoce & \textbf{Tardive} & Attention croisée\\
    \midrule
    Entraînement & oui & \textbf{non} & oui\\
    Modalité manquante & fragile & \textbf{robuste} & fragile\\
    Interprétable & non & \textbf{oui} & moyen\\
    \bottomrule
  \end{tabular}
  \medskip
  \[
    s_{\text{late}}(q,i)=\sum_m w_m\,(q\cdot z_i^m)=q\cdot\bar z_i,\qquad \bar z_i=\sum_m w_m z_i^m,\quad \sum_m w_m=1
  \]
  Un seul vecteur par item $\Rightarrow$ un seul index. $w$ réglé sur la validation (grille, rsum).
\end{frame}

% 8 -------------------------------------------------------------- Orateur : RAPHAËL
\begin{frame}{Audio : tête alignée par perte contrastive}
  \small
  Log-Mel (25\,ms, 10\,ms, 128 bandes) $\to$ AST (1\,212 tokens/10 s) $\to$ MLP $768\to1024\to512$ (1\,312\,256 paramètres).
  \[
    \mathcal L=\tfrac12\big(\mathcal L_{x\to y}+\mathcal L_{y\to x}\big),\quad
    \mathcal L_{x\to y}=-\tfrac1N\sum_i\log\tfrac{e^{S_{ii}}}{\sum_j e^{S_{ij}}},\quad S_{ij}=\tfrac{z^x_i\cdot z^y_j}{\tau}
  \]
  \[
    \frac{\partial\mathcal L_{x\to y}}{\partial z^x_i}=\frac1{N\tau}\sum_j(p_{ij}-\delta_{ij})\,z^y_j
  \]
  Ancre : légendes CLIP gelées ; une légende par clip et par lot ; $\tau$ appris (init. $0{,}07$, borné à $0{,}01$).
\end{frame}

\section{Application}
% 9 --------------------------------------------------------- Orateur : CLAUDY, RAPHAËL
\begin{frame}{Application Flask}
  \begin{columns}[T]
    \begin{column}{0.55\textwidth}
      \begin{itemize}
        \item Recherche texte, image ou audio (upload), top-$k$, scores cosinus et scores par modalité.
        \item Filtres : modalité, dataset, split, score minimal ; modes fusionné et mono-modal.
        \item Légendes d'un item, lecteur audio, classification zero-shot.
        \item Page embeddings : PCA, t-SNE, UMAP.
      \end{itemize}
    \end{column}
    \begin{column}{0.42\textwidth}
      \TODO{capture d'écran de l'application}
    \end{column}
  \end{columns}
\end{frame}

% 10 -------------------------------------------------------- Orateur : CLAUDY, RAPHAËL
\begin{frame}{Démonstration}
  \begin{enumerate}
    \item Texte, mode fusionné, puis mono-modal et filtre « image ».
    \item Image téléversée : légendes retrouvées.
    \item Audio téléversé, puis classification zero-shot.
    \item Page embeddings colorée par classe.
  \end{enumerate}
  \TODO{plan B : vidéo enregistrée}
\end{frame}

\section{Résultats}
% 11 --------------------------------------------------------------- Orateur : CLAUDY
\begin{frame}{Texte $\to$ image : mono-modal contre multimodal}
  \centering\figopt{fig_mono_vs_multi.png}{0.7\textheight}\\[2pt]
  {\small Flickr8k test, 1\,000 images, leave-one-out. \TODO{chiffres clés et gain}}
\end{frame}

% 12 ------------------------------------------------------------- Orateur : RAPHAËL
\begin{frame}{Audio : alignement et zero-shot ESC-50}
  \centering\figopt{fig_zero_shot_per_class.png}{0.7\textheight}\\[2pt]
  {\small \TODO{audio$\leftrightarrow$texte R@K ; top-1 et top-5 ESC-50 contre 2\,\% de hasard}}
\end{frame}

% 13 -------------------------------------------------------------- Orateur : FRESNEL
\begin{frame}{Ablations et espace commun}
  \centering
  \figopt{fig_dim_ablation.png}{0.5\textheight}\hspace{1em}\figopt{fig_proj_tsne.png}{0.5\textheight}\\[2pt]
  {\small \TODO{effet de $d$, poids $w$ retenu, alignement des modalités}}
\end{frame}

\section{Biais et conclusion}
% 14 ----------------------------------------------------------------- Orateur : IDA
\begin{frame}{Biais et limites de la fusion}
  \begin{itemize}
    \item Données : anglais, culture occidentale, personnes et animaux surreprésentés.
    \item Modalités : texte pivot ; image$\leftrightarrow$audio jamais appariés (évaluation par proxy).
    \item Alignement : écart entre modalités, scores non calibrés ; dérive Clotho $\to$ ESC-50.
    \item Fusion : poids global non adaptatif ; leave-one-out $\ne$ usage réel.
  \end{itemize}
  \centering\figopt{fig_bias_groups.png}{0.38\textheight}
\end{frame}

% 15 --------------------------------------------------------------- Orateur : CLAUDY
\begin{frame}{Conclusion et perspectives}
  \begin{itemize}
    \item \TODO{résultat principal chiffré, mono-modal contre multimodal}
    \item \TODO{résultat audio et zero-shot}
    \item Perspectives : fine-tuning partiel d'AST, CLAP, attention croisée, légendes en français, annotation humaine image--audio.
  \end{itemize}
  \centering\Large Merci. Questions ?
\end{frame}

\appendix
\begin{frame}{Annexe : exemple chiffré InfoNCE ($N=2$, $\tau=0{,}1$)}
  \small
  $x_1=(1,0),\ x_2=(0,1),\ y_1=(0{,}8660;0{,}5),\ y_2=(0{,}5;0{,}8660)$ \;$\Rightarrow$\; logits $8{,}660$ et $5{,}000$.\\
  $p_{11}=1/(1+e^{-3{,}66})=0{,}97491$, perte par ligne $=\ln1{,}0257325=0{,}025407$, donc $\mathcal L=0{,}025407$.\\
  Gradient sur $z^x_1$ : $\frac1{2\cdot0{,}1}\,0{,}025087\,(y_2-y_1)=(-0{,}04591;\,+0{,}04591)$.
\end{frame}

\begin{frame}{Annexe : budget de calcul}
  \small
  \begin{tabular}{lcc}
    \toprule
    Encodeur & MACs & Temps CPU (50 GFLOPS)\\
    \midrule
    CLIP texte & $2{,}98\cdot10^9$ & $\approx0{,}12$ s\\
    CLIP image & $4{,}29\cdot10^9$ & $\approx0{,}17$ s\\
    AST (10 s) & $1{,}30\cdot10^{11}$ & $\approx5{,}2$ s\\
    \bottomrule
  \end{tabular}
\end{frame}

\end{document}
